\textbf{结论名称：}裂项相消——对数型
\textbf{适用条件：}数列通项可写成 \(\log_a\!\bigl(1+\frac{1}{n}\bigr)\) 的形式，或能变形为该形式的对数差结构，且 \(a>0\)，\(a\neq1\)，\(n\in\mathbb{N}^*\)。
\textbf{变量说明：}底数 \(a>0\) 且 \(a\neq1\)；\(n\) 为正整数。
\textbf{核心公式：}
\[
   \boxed{\log_a\!\left(1+\frac{1}{n}\right)=\log_a(n+1)-\log_a n}
\]
\textbf{使用场景：}处理含对数与分式的数列求和时，将一项裂成两项之差，利用累加法使中间项相互抵消，简化求和。
\textbf{简单例子：}
求前 \(n\) 项和 \(S_n=\displaystyle\sum_{k=1}^{n}\log_2\!\left(1+\frac{1}{k}\right)\)。
裂项得 \(S_n=\sum_{k=1}^{n}\bigl[\log_2(k+1)-\log_2 k\bigr]=\log_2(n+1)-\log_2 1=\log_2(n+1)\)。
取 \(n=3\)，原式 \(=\log_2(2/1)+\log_2(3/2)+\log_2(4/3)=\log_2 4=2\)，公式给出 \(\log_2 4=2\)，一致。